\item \textbf{Problema de repaso: Colisión e integración con oscilador.}
Este problema extiende el razonamiento del problema 41 del capítulo 9. Dos deslizadores se ponen en movimiento sobre una pista de aire. El deslizador uno tiene masa $m_1 = 0.240\text{ kg}$ y se mueve hacia la derecha con una rapidez $0.740\text{ m/s}$. Tendrá una colisión posterior con el deslizador número dos, de masa $m_2 = 0.360\text{ kg}$, que tiene rapidez hacia la derecha $0.120\text{ m/s}$. Un resorte ligero con constante de fuerza de $45.0\text{ N/m}$ se une al extremo posterior del deslizador dos, como se muestra en la figura P9.41. Cuando el deslizador uno toca el resorte, un súper pegamento hace que instantánea e inmediatamente se pegue a su extremo del resorte.
\begin{enumerate}
    \item Encuentre la rapidez común que tienen los dos deslizadores cuando la compresión del resorte es un máximo.
    \item Encuentre la distancia máxima de compresión del resorte.
\end{enumerate}
El movimiento después de que los planeadores se unen consiste en una combinación de (1) el movimiento de velocidad constante del centro de masa del sistema de los dos planeadores encontrados en el inciso (1) y (2) movimiento armónico simple de los planeadores con respecto al centro de masa.
\begin{enumerate}
    \setcounter{enumi}{2}
    \item Encuentre la energía del movimiento de centro de masa.
    \item Encuentre la energía de oscilación.
\end{enumerate}